\section{Índice de Falar Sincero (IFS)}

O Índice de Falar Sincero (IFS) é uma medida projetada para avaliar o comportamento dos parlamentares em relação à disseminação de desinformação. A premissa central é quantificar a frequência e a intensidade com que um parlamentar propaga conteúdos semelhantes a narrativas comprovadamente falsas. A metodologia consiste em comparar as postagens de cada parlamentar com uma base vetorial de conteúdos previamente verificados como falsos por agências de checagem.

O processo de cálculo envolve as seguintes etapas:
\begin{enumerate}
    \item \textbf{Construção de um corpus de referência}: Criação de uma base vetorial a partir de conteúdos rotulados como falsos, que serve como referência para a detecção de desinformação.
    \item \textbf{Curadoria do corpus}: O material é filtrado para garantir a qualidade, com a remoção de duplicatas, ruídos e a seleção de alegações centrais.
    \item \textbf{Medição de similaridade}: As postagens de cada parlamentar são transformadas em embeddings e comparadas com os vetores do corpus de referência, utilizando a similaridade do cosseno.
    \item \textbf{Detecção de propagação}: Para cada postagem, identifica-se o grau de alinhamento com a desinformação. Uma postagem é considerada como potencialmente propagadora de conteúdo falso se sua similaridade máxima com alguma alegação do corpus ultrapassar um limiar $\tau$. Adicionalmente, um classificador de \textit{stance} é utilizado para distinguir se a postagem \textit{propaga} ou \textit{denuncia} o conteúdo.
\end{enumerate}

\subsection{Definição Operacional e Pontuação}

Para um parlamentar $i$, seja $\mathcal{S}_i$ o conjunto de pontuações de similaridade de suas postagens que foram classificadas como propagadoras de desinformação:
\[
    \mathcal{S}_i = \Big\{ s_{it} \;:\; s_{it} = \max_j \cos\big(e(\text{post}_{it}), e(\text{claim}_j)\big) \geq \tau \text{ e } \text{stance}_{it} = \text{propaga} \Big\}.
\]
Com base nesse conjunto, e sendo $P_i$ o total de postagens analisadas e $c_i = |\mathcal{S}_i|$ o número de postagens que propagam desinformação, duas formas de pontuação estão em avaliação. Ambas são orientadas de forma que, \textit{quanto maior o índice, melhor} o desempenho do parlamentar.

\begin{description}
    \item[A) Incidência (Taxa)] Pontua a ocorrência de propagação. O índice é maior quanto menor a taxa de postagens propagadoras.
    \[
        r_i = \frac{c_i}{P_i} \quad \Rightarrow \quad \mathrm{IFS}_i^{\text{taxa}} = 1 - r_i.
    \]

    \item[B) Distância (Similaridade Média)] Penaliza postagens que são vetorialmente mais próximas das alegações falsas.
    \[
        \bar{s}_i = \begin{cases}
            \frac{1}{c_i}\sum_{s \in \mathcal{S}_i} s, & \text{se } c_i > 0, \\
            0, & \text{se } c_i = 0
        \end{cases}
        \quad \Rightarrow \quad \mathrm{IFS}_i^{\text{dist}} = 1 - \bar{s}_i.
    \]
\end{description}

\subsection{Desafios Metodológicos}
A implementação do IFS enfrenta alguns desafios que exigem calibração e validação contínuas:
\begin{itemize}
    \item \textbf{Polaridade}: Distinguir com precisão entre \textit{denúncia} e \textit{propagação} de desinformação é um dos principais desafios, demandando modelos robustos de \textit{stance detection}.
    \item \textbf{Limiar de similaridade ($\tau$)}: A definição do limiar $\tau$ é um ponto crítico que afeta diretamente a sensibilidade e a especificidade do índice. Sua calibração deve ser realizada por meio de análises de curvas ROC/PR e validação com anotação humana para mitigar falsos positivos e negativos.
    \item \textbf{Cobertura do corpus}: A qualidade do índice depende da representatividade e da atualização contínua do corpus de referência, que deve abranger diferentes temas e períodos.
\end{itemize}

% \subsection{Validação e Aplicação}

Como parte da validação preliminar, foi realizada uma análise manual de postagens com base em checagens da agência Aos Fatos. O corpus de referência é populado a partir do arquivo \texttt{Analise manual Aos fatos\_v2.xlsx} e armazenado na coleção ChromaDB \texttt{checks\_v2\_rewrite\_cosine\_sentence\_transformer}. As tabelas abaixo apresentam os resultados agregados por parlamentar e partido.

\begin{table}[h!]
    \centering
    \caption{Resultados Preliminares do IFS por Parlamentar}
    \begin{tabular}{rrrr}
\toprule
Ano & IFS médio & Total de publicações & Quantidade de parlamentares \\
\midrule
2023 & 0.7 & 12 & 8 \\
2024 & 0.7 & 19 & 13 \\
\bottomrule
\end{tabular}

\end{table}

\begin{table}[h!]
    \centering
    \caption{Resultados Preliminares do IFS por Partido}
    \begin{tabular}{lrrr}
\toprule
Partido & Quantidade de parlamentares & Total de publicações & \% de parlamentares \\
\midrule
PL & 9 & 17 & 50.0 \\
PSOL & 2 & 3 & 11.1 \\
Avante  & 1 & 2 & 5.6 \\
Novo & 1 & 1 & 5.6 \\
MDB & 1 & 1 & 5.6 \\
PSB  & 1 & 1 & 5.6 \\
Podemos & 1 & 4 & 5.6 \\
Repúblicanos & 1 & 1 & 5.6 \\
União Brasil & 1 & 1 & 5.6 \\
\bottomrule
\end{tabular}

\end{table}

\subsection{Protocolo Experimental}

Embora métricas agregadas ainda não tenham sido reportadas, recomendamos o seguinte protocolo:

\begin{itemize}
    \item \textbf{Divisão dos dados:} 70/15/15 treino/validação/teste conforme definido em \texttt{config.py}.
    \item \textbf{Métricas:} Precisão, recall, F1-score, ROC-AUC e curvas de calibração (diagramas de confiabilidade) devido às saídas probabilísticas.
    \item \textbf{Ablações:} Comparar a cabeça MLP com uma cabeça linear; avaliar configurações com recuperação habilitada versus desabilitada; estudar perda focal contra entropia cruzada padrão.
    \item \textbf{Varreduras de hiperparâmetros:} Explorar taxas de dropout (0,3--0,6), larguras de camadas ocultas e parâmetros de perda focal.
\end{itemize}

\subsection{Análise da Distribuição}

A análise da distribuição do IFS, incluindo histogramas e boxplots, será detalhada nesta seção para ilustrar a dispersão e a concentração dos valores do índice entre os parlamentares.

\subsection{Notas de Implementação}

\begin{itemize}
    \item \textbf{Gerenciamento de dispositivo:} \texttt{config.DEVICE} escolhe automaticamente CUDA quando disponível.
    \item \textbf{Serialização:} Checkpoints podem armazenar um \texttt{state\_dict} bruto ou um dicionário contendo \texttt{model\_state\_dict}; ambos os formatos são suportados.
    \item \textbf{Reprodutibilidade:} Sementes são definidas via \texttt{config.RANDOM\_SEED}, mas determinismo de ponta a ponta requer semear NumPy, PyTorch CUDA e data loaders.
    \item \textbf{Dependências:} Pacotes-chave listados em \texttt{requirements.txt} incluem \texttt{transformers}, \texttt{torch}, \texttt{chromadb}, \texttt{pandas} e \texttt{sentence-transformers}.
\end{itemize}

\subsection{Limitações}

\begin{itemize}
    \item \textbf{Cobertura de dados:} O desempenho depende da cobertura de alegações verificadas e postagens anotadas; deriva de domínio pode prejudicar a precisão.
    \item \textbf{Latência:} Recuperação e inferência do transformador incorrem em sobrecarga computacional; processamento em lote e aceleração são essenciais para uso em tempo real.
    \item \textbf{Lacunas de avaliação:} A falta de benchmarks publicados limita a comparação com trabalhos anteriores.
\end{itemize}
